\chapter{2001年上海卷（秋文）}
\section{填空题}
\begin{problem}
设函数 \( f(x) = \log_a x \)，则满足 \( f(x) = \frac{1}{2} \) 的 \( x \) 值为 \fillinblank{} \fillinblank{}.
\end{problem}

\begin{problem}
设数列 \(\{a_{n}\}\) 的首项 \(a_1 = -7\)，且满足 \(a_{n+1} = a_n + 2 (n \in \mathbf{N})\)，则 \(a_1 + a_2 + \cdots + a_{17} = \fillinblank{}\).
\end{problem}

\begin{problem}
设 \(P\) 为双曲线 \(\frac{x^2}{4} - y^2 = 1\) 上一动点，\(O\) 为坐标原点，\(M\) 为线段 \(OP\) 的中点，则点 \(M\) 的轨迹方程为 \fillinblank{} \fillinblank{}.
\end{problem}

\begin{problem}
设集合 \(A = \{x \mid 2 \lg x = \lg (8x - 15), x \in \R\}\)，\(B = \{x \mid \cos \frac{\pi}{2} > 0 \in \R\}\)，则 \(A \cap B\) 的元素个数为 个 \fillinblank{}.
\end{problem}

\begin{problem}
抛物线 \(x^{2} - 4y - 3 = 0\) 的焦点坐标为 \fillinblank{} \fillinblank{}.
\end{problem}

\begin{problem}
设数列 \(\{a_{n}\}\) 是公比 \(q > 0\) 的等比数列，\(S_{n}\) 是它的前 \(n\) 项和. \(\lim_{n\to \infty}S_n = 7\)，则此数列的首项 \(a_1\) 的取值范围是 \fillinblank{} \fillinblank{}.
\end{problem}

\begin{problem}
某餐厅供应客食，每位顾客可以在餐厅提供的菜肴中任选 2 荤 2 素共 4 种不同的品种. 现在餐厅准备了 5 种不同的荤，若要保证每位顾客有 200 种以上不同的选择，则餐厅至少还需要准备不同的素菜品种 种. (结果用数值表示)
\end{problem}

\begin{problem}
在代数式 \(\left(x - \frac{1}{x^2}\right)^5\) 的展开式中，常数项为 \fillinblank{} \fillinblank{}.
\end{problem}

\begin{problem}
设 \(x = \sin \alpha, \alpha \in \left[-\frac{\pi}{6}, \frac{5\pi}{6}\right]\)，则 \(\arccos x\) 的取值范围为 \fillinblank{} \fillinblank{}.
\end{problem}

\begin{problem}
利用下列盈利表中的数据进行决策，应选择的方案是 \fillinblank{}. 自然状况 盈利(万元) 概率 方案 \( A_1 \) \( A_2 \) \( A_3 \) \( A_4 \) \( S_1 \) 0.25 50 70 -20 98 \( S_2 \) 0.30 65 26 52 82 \( S_3 \) 0.45 26 16 78 -10
\end{problem}

\begin{problem}
已知两个圆: \(x^{2} + y^{2} = 1\) \circled{1}与 \(x^{2} + (y - 3)^{2} = 1\) \circled{2}，则由\circled{1}式减去\circled{2}式可得上述两圆的对称轴方程. 将上述命题在曲线的情况下加以推广，即要求得到一个更一般的命题，而已知命题应成为所推广命题的一个特例，推广的命题为 \fillinblank{} \fillinblank{}.
\end{problem}

\begin{problem}
据报道，我国目前已成为世界上变荒漠化危害最严重的国家之一. 下左图表示我国土地沙化总面积在上个世纪五六十年代、七八十年代、九十年代的变化情况. 由图中的相关信息，可将上述有关年代中，我国年平均土地沙化面积在下右图中图示为:
\end{problem}

\section{单选题}
\begin{problem}
\(a = 3\) 是直线 \(ax + 2y + 3a = 0\) 和直线 \(3x + (a - 1)y = a - 7\) 平行且不重合的（\quad）

\choices
    {充分非必要条件}
    {必要非充分条件}
    {充要条件}
    {既非充分也非必要条件}
\end{problem}

\begin{problem}
在平行六面体 \(ABCD - A_{1}B_{1}C_{1}D_{1}\) 中，\(M\) 为 \(AC\) 与 \(BD\) 的交点，若 \(\overrightarrow{A_1B_1} = \overrightarrow{\alpha}, \overrightarrow{A_1D_1} = \overrightarrow{b}, \overrightarrow{A_1A} = \overrightarrow{c}\)，则下列向量中与 \(\overrightarrow{B_1M}\) 相等的向量是（\quad）

\choices
    {\(-\frac{1}{2}\overrightarrow{a} +\frac{1}{3}\overrightarrow{b} +\overrightarrow{c}\)}
    {\(\frac{1}{3}\overrightarrow{a} +\frac{1}{2}\overrightarrow{b} +\overrightarrow{c}\)}
    {\(\frac{1}{4}\overrightarrow{a} -\frac{1}{2}\overrightarrow{b} +\overrightarrow{c}\)}
    {\(-\frac{1}{2}\overrightarrow{a} -\frac{1}{2}\overrightarrow{b} +\overrightarrow{c}\)}
\end{problem}

\begin{problem}
已知 \(a, b\) 为两条不同的直线，\(\alpha, \beta\) 为两个不同的平面，且 \(a \perp \alpha, b \perp \beta\)，则下列命题中的假命题是（\quad）

\choices
    {若 \(a \nmid b\)，则 \(\alpha \nmid \beta\)}
    {若 \(\alpha \perp \beta\)，则 \(a \perp b\)}
    {若 \(a, b\) 相交，则 \(\alpha, \beta\) 相交}
    {若 \(\alpha, \beta\) 相交，则 \(a, b\) 相交}
\end{problem}

\begin{problem}
用计算器验算函数项 \( y = \frac{\lg x}{x} (x > 1) \) 的若干个值，可以猜想下列命题中的真命题只能是（\quad）

\choices
    {\(y = \frac{\lg x}{\lg x}\) 在 \((\mathrm{l}, + \infty)\) 上是单调减函数}
    {\( y = \frac{\lg x}{x}, x \in (1, +\infty) \) 的值域为 \( \left(0, \frac{\lg 3}{3}\right] \)}
    {\( y = \frac{\lg x}{\ln x}, x \in (1, +\infty) \) 有最小值}
    {\(\lim_{x\to \infty}\frac{\lg n}{n} = 0,n\in \mathbf{N}\)}
\end{problem}

\section{解答题}
\begin{problem}
已知 a、b、c 是 \( \triangle ABC \) 中 \( \angle A \) 、 \( \angle B \) 、 \( \angle C \) 的对边，S 是 \( \triangle ABC \) 的面积，若 a = 4, b = 5, \( S = 5\sqrt{3} \) ，求 c 的长度.
\end{problem}

\begin{problem}
设 \(F_{1}, F_{2}\) 为椭圆 \(\frac{x^{2}}{9} + \frac{y^{2}}{4} = 1\) 的两个焦点，\(P\) 为椭圆上的一点. 已知 \(P, F_{1}, F_{2}\) 是一个直角三角形的三个顶点，且 \(|PF_{1}| > |PF_{2}|\)，求 \(\frac{|PF_{1}|}{|PF_{2}|}\) 的值.
\end{problem}

\begin{problem}
在棱长为 \(a\) 的正方体 \(OABC - O'A'B'C'\) 中，\(E, F\) 分别是棱 \(AB, BC\) 上的动点，且 \(AE = BF\). (1) 求证: \(AF \perp C'E\); (2) 当三棱柱 \(B'' - BEF\) 的体积取得最大值时，求二面角 \(B'' - EF - B\) 的大小. (结果用反三角函数表示)
\end{problem}

\begin{problem}
对任意一个非零复数 \(z_{1}\) 定义集合 \(M_{z_1} = \{\omega | \omega = z^{2n - 1}, n \in \mathbf{N}\}\). (1) 设 \(a\) 是方程 \(x + \frac{1}{x} = \sqrt{2}\) 的一个根，试用列举法表示集合 \(M_{a_1}\). 若在 \(M_{z_1}\) 中任取两个数，求其和为零的概率 \(P\); (2) 设集合 \(M_{z_1}\) 中只有 3 个元素，试写出满足条件的一个 \(z\) 的值，并说明理由.
\end{problem}

\begin{problem}
用水清洗一堆蔬菜上残留的农药，对用一定量的水清洗一次的效果作如下假定：用1个单位量的水可洗掉蔬菜上残留农药用量的 \(\frac{1}{2}\)，用水越多次掉的农药量也越多，但总还有农药残留在蔬菜上. 设用 \(x\) 单位量的水清洗一次以后，蔬菜上残留的农药与本次清洗前残留在农药量之比为函数 \(f(x)\). (1) 试规定 \( f(0) \) 的值，并解释其实际意义； (2) 试根据假定写出函数 \( f(x) \) 应该满足的条件和具有的性质; (3) 设 \( f(x) = \frac{1}{1 + x^2 - 2x} \)，现有 \( \alpha (\alpha > 0) \) 单位量的水，可以清洗一次，也可以把水平均分成2份后清洗两次. 试问用哪种方案清洗后蔬菜上的农药量比较少？说明理由.
\end{problem}

\begin{problem}
对任意函数 \( f(x), x \in D \). 可按图示构造一个数列发生器，其工作原理如下: \circled{1} 输入数据 \(x_0 \in D\)，经数列发生器输出 \(x_1 = f(x_0)\); \circled{2} 若 \(x_{1} \notin D\)，则数列发生器结束工作；若 \(x_{1} \in D\)，则将 \(x_{1}\) 反馈回输入断，再输出 \(x_{2} = f(x_{1})\)，并依此规律继续下去. 规定式 \(f(x) = \frac{x - x_{1}}{x + 1}\). (1) 若输出 \(x_0 = \frac{49}{63}\)，则由数列发生器产生数列 \(\{x_n\}\)，请写出数列 \(\{x_m\}\) 的所有项; (2) 若要数列发生器产生一个无穷的常数数列，试求输出的初始数据 \(x_0\) 的值; (3) 是否存在 \(x_0\)，在输入数据 \(x_0\) 时，该数列发生器产生一个各项均为负数的无穷数列？若存在，求出 \(x_0\) 的值；若不存在，请说明理由.
\end{problem}

